\chapter{Discussão}\label{sec:discussao}

Este capítulo interpreta os resultados apresentados no \capref{cap4-resultados} à luz da questão de pesquisa inicial na \secref{sec-12-problema-de-pesquisa} e das taxonomias de paralelização discutidas na \secref{sec-26-gpu-metaheuristicas}.
A discussão está organizada em quatro partes: síntese dos achados e análise de memória; comparação com o estado da arte; limitações e ameaças à validade; e implicações práticas e diretrizes para o uso de GPUs na Pesquisa Operacional e na Engenharia.

\section{Síntese da pesquisa}\label{sec-51-sintese-achados}

A questão principal desta pesquisa investiga como comparar, de forma justa, o impacto de diferentes estratégias de paralelização em GPU sobre o desempenho de um algoritmo memético para o TSP (\secref{sec-12-problema-de-pesquisa}).
Os resultados do \capref{cap4-resultados} respondem a essa questão em duas frentes complementares: manter todo o ciclo evolutivo na VRAM produz os menores tempos em todas as \totalinstancias\ instâncias --- alcançando ganho médio geométrico de \(71{,}03\times\) sobre a CPU, frente a \(23{,}86\times\) da Híbrida Otimizada e \(4{,}57\times\) da Híbrida Ingênua (\tabref{tbl:results-aggregated}) ---, enquanto a qualidade das soluções não apresenta diferença estatisticamente significativa entre as quatro variantes sob os mesmos parâmetros de busca local (\(\kappa = \maxpassadas\) e \(k_{\text{viz}} = \kviz\), \secref{sec-43-analise-agregada}).
Na prática, portanto, a escolha entre as arquiteturas não envolve perda de acurácia, mas sim o equilíbrio entre velocidade de execução e consumo de memória na GPU.

O principal fator físico por trás das diferenças de tempo observadas na \tabref{tbl:phase-timings} e na \tabref{tbl:transfers-launches} é o contraste entre a largura de banda interna da VRAM (\(\approx 112\)\,GB/s) e a do barramento PCIe (\(\approx 16\)\,GB/s, \secref{sec-311-hardware-software}): mover dados pelo barramento custa cerca de sete vezes mais do que acessá-los dentro da placa, além da sobrecarga de cada lançamento de \emph{kernel}.
Por isso, na Híbrida Ingênua o excesso de travessias individuais desperdiça parte do ganho da GPU, enquanto a Híbrida Otimizada (\secref{sec-322-variantes}) reduz os lançamentos em mais de duas ordens de grandeza (\tabref{tbl:transfers-launches}) e passa a ser limitada pelo cruzamento em CPU (\secref{sec-44-decomposicao-fases}).
Por fim, a variante Totalmente em GPU elimina a transferência de matrizes de rotas por geração e concentra o processamento na VRAM; ainda assim, nas maiores instâncias avaliadas (como \texttt{pr1002}), em que a própria busca local em GPU volta a consumir a maior parte do tempo, o ganho adicional sobre a Híbrida Otimizada naturalmente se estreita.

Além do tempo de execução, o \textbf{comportamento de memória} forma um segundo eixo estrutural de análise, resumido na \tabref{tbl:memory-axis}.
No conjunto avaliado neste trabalho (até \(\maiorinstancia\) cidades, consumindo cerca de \(6\) a \(7\)\,MB na VRAM), todas as instâncias cabem com ampla folga nos \(4\)\,GB da placa (\secref{sec-311-hardware-software}).
Contudo, ao extrapolar para problemas muito maiores, a literatura de roteamento em GPU aponta que o termo quadrático da matriz densa (\(4\,n_{coords}^2\) bytes) rapidamente supera o termo linear da população (\(n_{pop} = 256\)) e passa a pressionar tanto a capacidade de VRAM quanto o uso de \emph{cache} e memória compartilhada \autocite{schulz2013gpu,oneil2015rethinking}.
Nesse regime, estratégias que mantêm dados residentes aceleram o acesso, mas exigem equilibrar desempenho e memória --- havendo alternativas dentro da própria GPU para contornar esse limite, ainda que geralmente menos eficientes em tempo.

\begin{table}[htbp]
      \centering
      \caption{Eixo de memória e comunicação por variante (análise estrutural, não medida experimental)}\label{tbl:memory-axis}
      \begin{tabelapadrao}
            \begin{tabular}{@{}
                        >{\raggedright\arraybackslash}p{(\linewidth - 4\tabcolsep) * \real{0.2200}}
                        >{\raggedright\arraybackslash}p{(\linewidth - 4\tabcolsep) * \real{0.4200}}
                        >{\raggedright\arraybackslash}p{(\linewidth - 4\tabcolsep) * \real{0.3600}}@{}}
                  \toprule
                  Variante                   & Perfil de transferência por geração                                                                                                                                                                           & Estado residente na GPU                                    \\
                  \midrule
                  \textbf{CPU}               & Nenhuma                                                                                                                                                                                                       & Nenhum                                                     \\
                  \sepitem
                  \textbf{Híbrida Ingênua}   & \(\mathcal{O}(n_{coords}^2 + n_{pop}\,n_{coords})\) por geração (reenvio de \(D\) e de cada circuito \(C_m\), ida e volta)                                                                                    & Matriz \(D\), lista \(N_k\) e circuito \(C_m\) em execução \\
                  \sepitem
                  \textbf{Híbrida Otimizada} & \(\mathcal{O}(n_{pop}\,n_{coords})\) por geração (\(P(t)\) enviada \(2\times\) \(\mathit{H2D}\) --- \(2\)-opt e avaliação --- e retornada \(1\times\) \(\mathit{D2H}\))                                       & Matriz \(D\), lista \(N_k\) e matriz \(P(t)\) na geração   \\
                  \sepitem
                  \textbf{Totalmente em GPU} & \(\mathcal{O}(n_{coords}^2 + n_{pop}\,n_{coords})\) na inicialização; \(\mathcal{O}(n_{pop})\) por geração (\(\approx 11{,}26\)\,KB de sorteios aleatórios \(\mathit{H2D}\); \(0{,}00\)\,MB \(\mathit{D2H}\)) & População \(P(t)\), matriz \(D\) e lista \(N_k\) completos \\
                  \bottomrule
            \end{tabular}
      \end{tabelapadrao}
\end{table}

Esse limite de memória, porém, pode ser contornado em problemas maiores por técnicas já conhecidas na literatura:

\begin{itemize}
      \item
            \textbf{Armazenamento triangular compactado:} explora a simetria do TSP para guardar apenas \(n_{coords}(n_{coords}-1)/2\) distâncias, reduzindo o espaço pela metade;
      \item
            \textbf{Recálculo de distâncias sob demanda:} evita materializar a matriz completa, trocando uso de memória por operações aritméticas durante a busca; isso, no entanto, tende a penalizar o tempo de execução;
      \item
            \textbf{Particionamento da vizinhança (\emph{neighbourhood splitting}):} divide os movimentos avaliados por vez quando a busca \(k\)-opt excede a memória disponível \autocite{schulz2013gpu};
      \item
            \textbf{Transferência em blocos:} envia fatias dos dados para a GPU quando o estado completo da instância não cabe simultaneamente na VRAM.
\end{itemize}

\section{Comparação com o estado da arte}\label{sec-52-comparacao-estado-arte}

Embora as variantes em GPU tenham alcançado alta velocidade de execução, é importante situar as limitações do AG+\(2\)-opt implementado frente ao estado da arte do TSP.
Algoritmos especializados, como o LKH \autocite{helsgaun2000effective} e o algoritmo genético baseado em \emph{Edge Assembly Crossover} (EAX) \autocite{nagata2013powerful}, atingem com frequência \(0\%\) de \emph{gap} (solução ótima) na maioria das instâncias padrão da TSPLIB.

O AG+\(2\)-opt avaliado neste trabalho adota operadores clássicos justamente para isolar, de forma controlada, o impacto de cada estratégia de paralelização em GPU.
O foco, portanto, está na comparação justa entre variantes do mesmo algoritmo, e não na disputa por poder máximo de refinamento contra resolvedores especializados que empregam buscas locais de profundidade variável e operadores avançados de reconexão de arestas.
Trabalhos futuros podem investigar a adaptação desses operadores mais sofisticados (como os movimentos de Lin--Kernighan ou o EAX) para GPU, reduzindo essa diferença de qualidade.

Essa delimitação ajuda a ler corretamente os resultados do \capref{cap4-resultados}: os \emph{gaps} obtidos (de \(0{,}04\%\) em \texttt{berlin52} a \(1{,}74\%\) em \texttt{pr1002}, com média global em torno de \(0{,}90\%\)) não medem competitividade absoluta frente a resolvedores especializados --- sobretudo porque, nas instâncias menores, o critério de parada por alvo (\(\text{gap} \le \varepsilon_{\text{alvo}} = 1\%\)) interrompe deliberadamente a evolução assim que esse limiar é cruzado ---, mas confirmam que a paralelização preserva a capacidade de busca do algoritmo base.

\section{Possíveis limitações e ameaças à validade}\label{sec-53-ameacas-validade}

\subsection{Controle de parâmetros e critério de parada}\label{sec-531-controle-parametros}

Em estudos comparativos de meta-heurísticas paralelas, um risco frequente é comparar implementações com orçamentos computacionais diferentes.
Neste trabalho, todas as variantes usaram exatamente os mesmos parâmetros de busca local: limite de \(\kappa = \maxpassadas\) varreduras \(2\)-opt por indivíduo e lista estática com \(k_{\text{viz}} = \kviz\) vizinhos mais próximos (\secref{sec-321-ga-base}).
Isso garante que as diferenças de tempo observadas decorram da estratégia de paralelização e do fluxo de dados, e não de esforço algorítmico desigual.
Cabe registrar, porém, duas ressalvas metodológicas ligadas aos critérios de parada (\tabref{tbl:ga-hyperparameters}): primeiro, os limites por estagnação e por número máximo de gerações escalam com \(n_{coords}\), concedendo orçamento maior às instâncias maiores; segundo, o critério primário \(\text{gap} \le \varepsilon_{\text{alvo}} = 1\%\) é \emph{oracular}, pois depende do conhecimento prévio do ótimo \(B^*\) da TSPLIB.
Na prática, esse desenho caracteriza uma avaliação do tipo tempo até o alvo (\hyperlink{term:time-to-target}{\emph{time-to-target}}), metodologicamente válida para comparar velocidade até uma meta fixa de qualidade em \emph{benchmarks}, mas distinta de cenários industriais reais nos quais \(B^*\) é desconhecido e a execução depende apenas de estagnação ou de limite de tempo.

\subsection{Diferenças de execução no operador \texorpdfstring{\(2\)-opt}{2-opt}}\label{sec-5311-semantica-operadores}

Mesmo com parâmetros idênticos, a forma de avaliar a vizinhança muda entre CPU e GPU.
Em todas as variantes, cada cidade consulta os mesmos \(k_{\text{viz}} = \kviz\) vizinhos mais próximos.
Na CPU, os candidatos são percorridos em um laço vetorizado que seleciona o melhor movimento da varredura; na GPU, os mesmos candidatos são distribuídos entre as \emph{threads} de um bloco, que elegem a melhor troca por redução paralela.

Como a fórmula inteira de ganho e a lista de candidatos são exatamente iguais, a Híbrida Ingênua e a Híbrida Otimizada (\secref{sec-322-variantes}) produzem custos finais idênticos bit a bit em todas as \totalinstancias\ instâncias e em todas as \totalrepeticoes\ sementes avaliadas, confirmando que a diferença no modo de envio apenas reduz a sobrecarga de comunicação sem alterar o trajeto da busca.
Já entre CPU e GPU, a correspondência não é bit a bit: quando dois movimentos têm exatamente o mesmo ganho, a ordem de desempate em reduções paralelas na GPU pode diferir da varredura sequencial na CPU (\secref{sec-32-arquitetura-variantes}), sem provocar diferença estatisticamente significativa na qualidade final das soluções (\secref{sec-43-analise-agregada}).

\subsection{Limitações na medição de desempenho}\label{sec-532-instrumentacao}

As medições coletadas cobrem tempo total, tempo por etapa no relógio da CPU, volume transferido pelo barramento PCIe e número de lançamentos de \emph{kernel}, mas não incluem consumo de energia, evolução dinâmica da memória durante a execução nem contadores internos de baixo nível da GPU.
Além disso, pequenas parcelas de tempo fora das etapas cronometradas e o custo fixo de lançamento de \emph{kernels} são difíceis de isolar sem introduzir sobrecarga adicional na própria medição (\secref{apendice-telemetria}).
Como mostram os resultados da \tabref{tbl:phase-timings}, essa parcela residual tem peso maior em instâncias menores --- em que a execução dura poucas dezenas de milissegundos ---, mas tende a se tornar desprezível conforme o número de cidades aumenta (com a soma das etapas fechando entre \(99{,}8\%\) e \(100{,}0\%\) nas maiores instâncias avaliadas).

Duas limitações instrumentais específicas dessa telemetria merecem destaque explícito.
Primeiro, na variante \textbf{Totalmente em GPU}, os \emph{kernels} são disparados de modo assíncrono pela CPU e ela só aguarda a conclusão de todos eles no ponto de sincronização ao final da geração, fazendo com que até \(94{,}8\%\) do tempo medido pelo relógio da CPU se acumule na etapa \texttt{Sinc.} em \texttt{pr1002} (\tabref{tbl:phase-timings}); essa característica impede decompor internamente quanto tempo o dispositivo dedicou à busca local \(2\)-opt frente ao cruzamento OX sem forçar uma pausa de espera após cada operador (o que distorceria o tempo total coletado).
Segundo, na coleta experimental das variantes Híbrida Otimizada e Totalmente em GPU, o contador em software da telemetria somou \(8\,n_{coords}^2\) bytes adicionais (\(8{,}03\)\,MB em \texttt{pr1002}) pela recontagem do ponteiro compartilhado da matriz \(D\) em VRAM entre módulos independentes, sem transferência física correspondente no barramento PCIe --- artefato inteiramente reconciliado na \secref{sec-44-decomposicao-fases} e na nota da \tabref{tbl:transfers-launches}.

\subsection{Dependência do hardware utilizado}\label{sec-533-hardware-unico}

Todos os experimentos foram executados em um único equipamento, combinando um processador \textbf{Intel Core i7-7700HQ} e uma GPU de entrada \textbf{NVIDIA GeForce GTX 1050 Mobile} com \(4\)\,GB de VRAM (\secref{sec-311-hardware-software}).
Como o ganho de velocidade expressa a razão entre os tempos da CPU e da GPU, os valores numéricos exatos dependem diretamente do equilíbrio entre os dois componentes --- e não apenas da placa gráfica.

Do lado da \textbf{CPU}, processadores mais recentes, com mais núcleos, \emph{caches} maiores, frequências mais altas ou instruções vetoriais mais amplas --- ou ainda uma implementação em CPU paralelizada em múltiplas \emph{threads} --- tendem a reduzir o tempo de referência e diminuir o ganho relativo da GPU.
Do lado da \textbf{GPU}, modelos com mais unidades de processamento, maior capacidade de VRAM e maior largura de banda interna ampliam o porte das instâncias comportadas na memória e deslocam o ponto em que a comunicação ou a busca local passa a dominar o tempo.
Além disso, diferenças de limite térmico e de energia entre \emph{notebooks}, computadores de mesa e servidores também influenciam os tempos medidos.
Assim, os fatores exatos de aceleração valem para o equipamento avaliado, enquanto a ordenação geral entre as estratégias de paralelização tende a se manter em outros cenários de \emph{hardware}.

\section{Implicações práticas e aplicações na Engenharia}\label{sec-54-implicacoes-praticas}

Além dos resultados específicos sobre o TSP, a experiência experimental deste trabalho --- interpretada à luz da literatura de programação paralela \autocite{kirk2010programming,schulz2013gpu,oneil2015rethinking,nvidia2024cuda} --- oferece diretrizes diretas para o projeto de meta-heurísticas em GPU na Pesquisa Operacional e em problemas de Engenharia Mecânica e Industrial.
Nessas áreas, a mesma estrutura de sequenciamento e roteamento aparece no planejamento de trajetórias de ferramentas em máquinas CNC (como furação de chapas e corte a laser, reduzindo deslocamentos em vazio), no sequenciamento de pontos de solda e inspeção por robôs industriais e na logística interna e externa de fábricas; acelerar o ciclo de busca permite explorar mais cenários de projeto ou replanejar operações em tempo próximo do real.

Sob essa perspectiva, cinco diretrizes práticas sintetizam as lições dos experimentos (lembrando que as quatro variantes comparadas ilustram pontos representativos do espaço de projeto):

\begin{enumerate}
      \item
            \textbf{Manter o estado na VRAM e estruturar \emph{kernels} adequados à hierarquia de memória:} Evitar transferências individuais por rota (\textbf{Híbrida Ingênua}) e concentrar os dados na VRAM (\textbf{Totalmente em GPU}) é decisivo, mas insuficiente sem \emph{kernels} bem projetados --- distribuição equilibrada de \emph{threads}, acesso coalescido, memória compartilhada e \hyperlink{term:bitmap}{máscaras de bits} \autocite{kirk2010programming,nvidia2024cuda}. Quando problemas muito maiores excederem a VRAM ou a memória compartilhada, as estratégias de contorno elencadas na \secref{sec-51-sintese-achados} (armazenamento compacto, cálculo sob demanda, fatiamento em janelas ou arquiteturas híbridas agrupadas com \emph{streams}) preservam a viabilidade com penalidade controlada de tempo \autocite{schulz2013gpu,oneil2015rethinking,tsp_gpu}.
      \item
            \textbf{Eliminar gargalos remanescentes em CPU à luz da Lei de Amdahl:} Depois que a busca local é acelerada na GPU, etapas aparentemente leves mantidas na CPU --- como o cruzamento OX na Híbrida Otimizada (\secref{sec-44-decomposicao-fases}) --- passam a dominar o tempo total conforme a Equação~\eqref{eq:amdahl}. O ganho prático em algoritmos iterativos depende do equilíbrio de todo o ciclo evolutivo.
      \item
            \textbf{Reformular desvios condicionais para o modelo SIMT:} Como os \emph{warps} da GPU executam instruções em passo coordenado (\secref{sec-263-cpu-vs-gpu}), ramos condicionais irregulares serializam a execução e reduzem a ocupação dos núcleos \autocite{kirk2010programming,nvidia2024cuda}. Aproveitar o acelerador requer traduzir decisões para operações matriciais, reduções paralelas ou mapas de bits, em vez de portar diretamente a lógica sequencial da CPU.
      \item
            \textbf{Explorar a estrutura geométrica do problema em ambas as plataformas:} Restringir a busca local aos \(k_{\text{viz}} = \kviz\) vizinhos mais próximos reduziu expressivamente o esforço computacional sem degradação observada de qualidade. Em problemas geométricos de engenharia --- como furos adjacentes em uma chapa ou estações próximas em uma célula fabril ---, listas de candidatos ampliam o alcance do \emph{hardware} disponível tanto em CPU quanto em GPU.
      \item
            \textbf{Adotar medição controlada e validação estatística rigorosa:} Inicializar e compilar os \emph{kernels} (\hyperlink{term:warmup}{\emph{warmup}}) fora da janela cronometrada, fixar sementes aleatórias e aplicar protocolos não paramétricos (\secref{sec-34-protocolo-estatistico}) assegura que escolhas arquiteturais e investimentos em aceleradores se apoiem em evidências reproduzíveis.
\end{enumerate}
